Derivación propia, escrita antes de ejecutar el motor. No se copia una tabla ni una ecuación numerada de una norma. Si una norma de acero escribe un amplificador de piso parecido, la página, la sección y el número de ecuación no están en esta sesión: cita desconocida. La fórmula que se acepta es la de este apartado.

\paragraph{Qué aproximación es \(B_2\).}
Se llama \(B_2\) al cociente entre la deriva de piso de segundo orden y la deriva de piso de primer orden, cuando el único efecto geométrico que se retiene es el de la cuerda: la carga vertical \(\Sigma P\), al desplazarse \(\Delta\), equivale a un cortante adicional \((\Sigma P)\,\Delta/h\).

Sea un piso de altura \(h\), cortante de piso \(H\) y rigidez lateral \(K\) medida en el análisis lineal, de modo que la deriva de primer orden es \(\Delta_1=H/K\). El equilibrio linealizado del piso, con la axial de gravedad constante, es
\begin{equation}
K\Delta = H + \frac{\Sigma P}{h}\Delta.
\label{eq:val037-equilibrio}
\end{equation}
De ahí
\begin{equation}
\Delta = \frac{\Delta_1}{1-\Sigma P/(Kh)},\qquad
P_{e,\mathrm{piso}}=Kh=\frac{H\,h}{\Delta_1},
\label{eq:val037-pe}
\end{equation}
y el amplificador de la forma cerrada es
\begin{equation}
B_2^{\mathrm{piso}}=\frac{1}{1-\Sigma P/P_{e,\mathrm{piso}}}.
\label{eq:val037-b2}
\end{equation}
\(\Sigma P\) es la compresión del piso (positiva) y \(P_{e,\mathrm{piso}}\) es una sola carga crítica de piso, obtenida de su rigidez lateral. La suma del enunciado es esa carga: no se reparte otra vez entre columnas. La fórmula supone deriva única del piso, axial que no cambia con el corrimiento lateral, y ausencia del arqueo dentro de la barra. No es un factor de norma ni lleva un coeficiente de mayoración.

Esa última hipótesis es la que el motor no cumple de forma exacta. El método \texttt{pdelta} congela \(K_g\) en el desplazamiento de la gravedad: es pretensión inicial, la misma hipótesis de axial constante. La \(K_g\) de la barra, ya derivada en \texttt{provenance/ENG-ELM-002.md}, es la segunda variación \(N\int (v')^2\,\mathrm{d}x\) con las formas cúbicas de Hermite y \(N\) positiva en tracción:
\[
K_{gb}=\frac{N}{30L}
\begin{bmatrix}
36 & 3L & -36 & 3L\\
3L & 4L^2 & -3L & -L^2\\
-36 & -3L & 36 & -3L\\
3L & -L^2 & -3L & 4L^2
\end{bmatrix}.
\]
En un corrimiento de piso con giros impedidos, el vector transversal es \((0,0,\Delta,0)\). La entrada de cortante vale \(36\,N/(30h)=(6/5)\,N/h\). Con compresión \(N=-P\),
\begin{equation}
B_2^{\mathrm{Hermite}}=\frac{1}{1-1{,}2\,\Sigma P/P_{e,\mathrm{piso}}}.
\label{eq:val037-hermite}
\end{equation}
El \(6/5\) es el arco de la cúbica: \(\int_0^h (v')^2\,\mathrm{d}x=(6/5)\Delta^2/h\), frente a \(\Delta^2/h\) de la cuerda recta. Por eso \(B_2\) de la ecuación~\eqref{eq:val037-b2} es una aproximación del operador del motor, no su solución exacta. Partir la columna en más elementos acercaría el resultado a la viga-columna \(EI\,w''''+P\,w''=0\) con giros nulos en los extremos, cuya expansión también empieza por \(1+1{,}2\,r\) y no por \(1+r\). El hueco no se cierra refinando la malla.

Con \(r=\Sigma P/P_{e,\mathrm{piso}}\) y \(r\) pequeño,
\[
B_2^{\mathrm{piso}}=1+r+O(r^2),\qquad
B_2^{\mathrm{Hermite}}=1+1{,}2\,r+O(r^2).
\]
La diferencia es del orden de \(0{,}2\,(B_2^{\mathrm{piso}}-1)\).

\paragraph{Modelo que se va a ejecutar.}
Pórtico plano de tres pisos y un vano, en el plano \(XY\), \(Y\) hacia arriba. Cada piso mide \(h=4~\mathrm{m}\) y el vano \(L=6~\mathrm{m}\). Hay un elemento de Euler--Bernoulli por barra, sin relajaciones. \(E=2\cdot 10^{8}~\mathrm{Pa}\). Columnas y vigas tienen \(A=0{,}2~\mathrm{m}^2\). La inercia de cada columna es \(I_c=0{,}02~\mathrm{m}^4\) y la de cada viga \(I_v=6~\mathrm{m}^4\).

El cociente de rigideces de nudo es \(G=\sum(I_c/h)/\sum(I_v/L)\). Aquí \(I_c/h=0{,}005~\mathrm{m}^3\) e \(I_v/L=1~\mathrm{m}^3\). En la cubierta, un solo pilar, \(G=0{,}005\). En un piso intermedio, dos pilares, \(G=0{,}01\). Los nudos casi no giran: el error que eso añade a \(B_2\) es del orden de \(G\,r\), menor que \(10^{-3}\) en el régimen de abajo.

Cada piso recibe \(1{,}6\cdot 10^{5}~\mathrm{N}\) hacia \(-Y\), repartidos por igual en los dos nudos, y \(10^{3}~\mathrm{N}\) hacia \(+X\), también repartidos. \(\Sigma P\) de un piso es la suma de esas cargas verticales en ese piso y en los de arriba. \(H\) es la suma de las cargas \(+X\) en ese piso y en los de arriba. No se leen axiles recuperados: con cargas nodales verticales y simetría, la viga no necesita axil para el equilibrio vertical y las dos columnas se reparten \(\Sigma P\).

Si los giros fueran nulos, cada columna aportaría \(12EI_c/h^3=7{,}5\cdot 10^{5}~\mathrm{N/m}\) y el piso \(K=1{,}5\cdot 10^{6}~\mathrm{N/m}\), así \(P_{e,\mathrm{piso}}=Kh=6{,}0\cdot 10^{6}~\mathrm{N}\). Los tres cocientes de diseño, de cubierta a base, serían \(r=0{,}0267\), \(0{,}0533\) y \(0{,}080\). Son una cota de rigidez: con vigas finitas el piso es un poco más flexible, \(r\) sube un poco y no pasa de ese orden. Esta cota no es el valor de aceptación. \(P_{e,\mathrm{piso}}\) de la comparación sale de las derivas lineales del propio motor, \(H\,h/\Delta_1\).

\(B_2\) del motor es \(\Delta_2/\Delta_1\). \(\Delta_1\) es la deriva de piso del caso lateral en solitario, con \texttt{linearStatic}. \(\Delta_2\) es la deriva del mismo caso lateral con \texttt{pdelta}, cuya gravedad es el caso vertical. La media de \(U_x\) de los dos nudos del piso, menos la del piso de abajo, define la deriva. En la base esa media es nula.

El control de la iteración de Picard, distinto de la banda de aceptación, queda así: tolerancia relativa \(10^{-10}\), cuarenta pasos como máximo, escala de desplazamiento \(10^{-3}~\mathrm{m}\) y escala de giro \(10^{-3}~\mathrm{rad}\). Con \(r\) del orden de \(0{,}1\), el error de la iteración queda muy por debajo de \(10^{-8}\) en el desplazamiento y no mueve \(B_2\) dentro de la banda.

\paragraph{Tolerancia preregistrada.}
Se acepta el piso cuando
\begin{equation}
\bigl|B_2^{\mathrm{motor}}-B_2^{\mathrm{piso}}\bigr|
\le 0{,}35\,\bigl(B_2^{\mathrm{piso}}-1\bigr)+0{,}005.
\label{eq:val037-tol}
\end{equation}
El coeficiente \(0{,}35\) cubre el \(0{,}2\) del arqueo cúbico y el término \(O(r^2)\), con margen, mientras \(r<0{,}25\). En \(r=0{,}25\) el cociente entre el hueco de la ecuación~\eqref{eq:val037-hermite} y \((B_2^{\mathrm{piso}}-1)\) sigue por debajo de \(0{,}30\). El sumando \(0{,}005\) cubre el acoplamiento de los tres pisos con \(G\le 0{,}01\) y el redondeo de \texttt{f64}; es mayor que \(G\,r\) y mucho mayor que la tolerancia de Picard.

La banda no da por bueno un motor sin \(K_g\). Ese motor daría \(B_2=1\), a una distancia \(B_2^{\mathrm{piso}}-1\) de la forma cerrada. Esa distancia queda fuera de~\eqref{eq:val037-tol} en cuanto \(B_2^{\mathrm{piso}}-1>0{,}005/0{,}65\approx 0{,}0077\). Por eso cada piso debe cumplir, además, \(B_2^{\mathrm{piso}}-1>0{,}02\) y \(0<\Sigma P/P_{e,\mathrm{piso}}<0{,}25\). Si un piso sale de ese régimen, la comparación no se declara cumplida: no se ensancha la banda.

Con la cota de nudos rígidos, el piso bajo da \(B_2^{\mathrm{piso}}=1/0{,}92\approx 1{,}087\) y \(B_2^{\mathrm{Hermite}}\approx 1{,}106\). El hueco es aproximadamente \(0{,}019\) y la banda aproximadamente \(0{,}035\). En la cubierta el hueco es aproximadamente \(0{,}006\) y la banda aproximadamente \(0{,}015\). Esos dos números son aritmética de la cota, no una salida del motor.

\paragraph{Estado de los métodos.}
\texttt{buckling.linear} es Experimental en el registro de la fachada. \texttt{pdelta} (fachada \texttt{pDelta}) está publicado como Stable en ese mismo registro. Este caso no cambia el registro ni el código de producción. La comparación de \(B_2\) es analítica y no es una conformidad de norma. Se enuncia el estado de los dos métodos porque el encargo lo pide: la verificación de este piso trata el amplificador como una aproximación todavía abierta, con la banda de arriba, aunque la fachada no marca \texttt{pDelta} como Experimental.
